\documentclass[11pt]{article}
\usepackage[margin=1in]{geometry}
\usepackage{amsmath,amssymb,amsthm}
\usepackage{booktabs}
\usepackage{longtable}
\usepackage{hyperref}

\newtheorem{proposition}{Proposition}
\newtheorem{theorem}{Theorem}
\newtheorem{corollary}{Corollary}
\newtheorem{lemma}{Lemma}
\theoremstyle{remark}
\newtheorem{remark}{Remark}

\newcommand{\E}{\mathbb{E}}
\newcommand{\Var}{\operatorname{Var}}
\newcommand{\dd}{\,\mathrm{d}}

\title{Asymmetric Capital Buffer Control}
\author{}
\date{9 August 2026}

\begin{document}
\maketitle

\begin{abstract}
\noindent
A financial intermediary manages a capital buffer measured against a
reference level, penalised asymmetrically for shortfall. This document
states what the model establishes. Every numbered result carries a
canonical verifier, recorded in Appendix~\ref{app:status}: a Lean proof
where one exists, an exact symbolic check otherwise. The numerical
findings of Section~\ref{sec:numerical} are labelled as such and are
findings at stated parameters, not theorems. Section~\ref{sec:open}
records what is not established.
\end{abstract}

\tableofcontents

\section*{Index of result codes}
\addcontentsline{toc}{section}{Index of result codes}
\noindent
Every numbered result carries a three-letter code, used for reference in
the accompanying summaries. Codes are stable; numbering may change.

\begin{longtable}{lll}
\toprule
Code & Result & Statement \\
\midrule
\endhead
ACI & Prop.~\ref{prop:identity} & Accounting identity \\
PSA & Prop.~\ref{prop:persistence} & Persistence asymmetry \\
FPV & Thm.~\ref{thm:lambda4} & Fourth-power variance ratio \\
ROB & Prop.~\ref{prop:robust} & Robustness of the variance ratio \\
EGE & Prop.~\ref{prop:egarch} & EGARCH-class equivalence \\
SJT & Prop.~\ref{prop:sojourn} & Sojourn time \\
NST & Prop.~\ref{prop:nesting} & Exact nesting of the risk-neutral benchmark \\
RPC & Prop.~\ref{prop:concave-payoff} & The running payoff is concave \\
CCP & Prop.~\ref{prop:rrL} & Consistency condition on parameters \\
LQB & Prop.~\ref{prop:lowerbound} & Lower bound, the liquidation payoff \\
TSO & Prop.~\ref{prop:statemap} & The two states are one state \\
CDA & Cor.~\ref{cor:boundary} & The degeneracy is a coordinate artefact \\
SCG & Prop.~\ref{prop:satlimits} & The surplus limit is cap geometry too \\
REG & Lem.~\ref{lem:reg} & Regularity in the parameters \\
ENV & Lem.~\ref{lem:envelope} & The envelope identity \\
BDR & Lem.~\ref{lem:bdr} & Differentiability of the boundaries \\
TCS & Prop.~\ref{prop:tcs} & Threshold comparative statics \\
TCSN & Rem.~\ref{rem:tcsn} & What the proposition rests on \\
ENVK & Lem.~\ref{lem:envelopeK} & The envelope identity for the fixed cost \\
KCS & Prop.~\ref{prop:kcs} & Fixed-cost comparative statics \\
IDN & Cor.~\ref{cor:idn} & Joint local identification \\
KCSN & Rem.~\ref{rem:kcsn} & Reading the identification result \\
SOC & Prop.~\ref{prop:soc} & Second-order conditions \\
SMF & Prop.~\ref{prop:smf} & Smooth fit at the trigger \\
SMFN & Rem.~\ref{rem:smfn} & Why the argument discriminates \\
SOCN & Rem.~\ref{rem:socn} & Strictness at the trigger \\
ITR & Lem.~\ref{lem:itr} & Interior regularity \\
FIN & Lem.~\ref{lem:fin} & No chattering \\
VER & Prop.~\ref{prop:ver} & Verification \\
VERN & Rem.~\ref{rem:vern} & What verification needs \\
SLV & Rem.~\ref{rem:solver} & The solver, and how it was validated \\
CBI & Rem.~\ref{rem:capbinds} & The cap binds throughout the inaction region \\
CSL & Rem.~\ref{rem:compstat} & Comparative statics in the asymmetry parameter \\
LNC & Rem.~\ref{rem:nonconcave} & Local non-concavity under a fixed issuance cost \\
KSW & Rem.~\ref{rem:Ksens} & The fixed cost is the cause, a $K$-sweep \\
\bottomrule
\end{longtable}


\bigskip
\noindent\textbf{Terminology, binding throughout.} The parameter
$\lambda \ge 1$ is an \emph{asymmetric risk-response parameter}. It is not
called loss aversion. The distinction matters because $\lambda = 1$ recovers
a risk-neutral benchmark \emph{exactly} (Proposition~\ref{prop:nesting}), so
$\lambda$ measures asymmetry in excess of what optimisation under regulatory
constraint already produces --- not in excess of zero. Loss-aversion
language would additionally require elicited-choice evidence of the kind
observational data cannot supply.

\medskip
\noindent\textbf{Two parameters share the symbol $\lambda$.} The
discrete-time variance ratio of Section~\ref{sec:discrete} is written
$\lambda_V$; the shortfall penalty of Section~\ref{sec:formulation} is
written $\lambda_S$. They are distinct objects on the same state
(Proposition~\ref{prop:statemap}), and no relation between them is
asserted here; see Section~\ref{sec:open}.

\section{The discrete-time model}
\label{sec:discrete}

\subsection{Setup}

Let $E_t$ denote book equity and $Q_t$ required capital, with $Q_t$
\emph{exogenous} and growing deterministically,
$Q_t = (1+g)Q_{t-1}$, $g > -1$. The state is the log capital-buffer ratio
\begin{equation}
z_t \;=\; \log\!\left(\frac{E_t}{Q_t}\right),
\label{eq:z}
\end{equation}
so $z_t = 0$ exactly when equity sits at its requirement; $z_t > 0$ is
surplus and $z_t < 0$ is deficit. Equity accumulates by retention,
$E_t = E_{t-1}(1 + r_E b)$ where $r_E$ is the return on equity and
$b \in (0,1]$ the retention ratio.

Shocks enter with a state-dependent conditional variance,
\begin{equation}
z_t = \phi\, z_{t-1} + \varepsilon_t, \qquad
\Var(\varepsilon_t \mid z_{t-1}) = V(z_{t-1})
= V_0 \exp\!\big(-2(\log\lambda_V) \tanh(z_{t-1}/\theta)\big),
\label{eq:zrec}
\end{equation}
with $\theta>0$ a saturation scale. The mechanism is reduced-form: trading
limits tighten as the buffer erodes, so the same shock has a larger
conditional variance in deficit than in surplus. Nothing in
\eqref{eq:zrec} is derived from an optimisation; that is the role of
Section~\ref{sec:formulation} onward.

\subsection{Results}

\begin{proposition}[ACI: Accounting identity]
\label{prop:identity}
Balanced growth of equity and requirement forces $r_E = g/b$.
\end{proposition}

\begin{proof}
From $E_t = E_{t-1}(1+r_E b)$ and $Q_t = Q_{t-1}(1+g)$,
\begin{equation}
\frac{E_t}{Q_t} \;=\; \frac{E_{t-1}}{Q_{t-1}} \cdot \frac{1+r_E b}{1+g}.
\label{eq:prop1-ratio}
\end{equation}
Balanced growth means $E_t/Q_t = E_{t-1}/Q_{t-1}$ for every $t$ along the
deterministic path, i.e.\ the factor multiplying $E_{t-1}/Q_{t-1}$ in
\eqref{eq:prop1-ratio} equals $1$:
\begin{equation}
\frac{1+r_E b}{1+g} = 1
\;\Longleftrightarrow\;
1+r_E b = 1+g
\;\Longleftrightarrow\;
r_E b = g
\;\Longleftrightarrow\;
r_E = \frac{g}{b}.
\label{eq:prop1-chain}
\end{equation}
Each step of \eqref{eq:prop1-chain} is reversible, so this is not merely
necessary but exactly characterises balanced growth: given $b\in(0,1]$,
$r_E=g/b$ is well-defined for every $g>-1$ and is the unique return on
equity for which \eqref{eq:prop1-ratio} holds with equality at every $t$.
\end{proof}

\begin{proposition}[PSA: Persistence asymmetry]
\label{prop:persistence}
The deficit-side persistence is
$\phi^- = 1 - g(1-b)/b$, and $\phi^- \to 1$ exactly as $g \to 0$. The
surplus-side persistence $\phi^+$ differs, giving an asymmetry
$\phi^- - \phi^+$ which vanishes with $g$.
\end{proposition}

\begin{theorem}[FPV: Fourth-power variance ratio]
\label{thm:lambda4}
Under \eqref{eq:zrec}, the ratio of conditional variances deep in deficit
to deep in surplus is
\begin{equation}
\frac{\lim_{z\to-\infty} V(z)}{\lim_{z\to+\infty} V(z)}
\;=\; \frac{V_0\lambda_V^2}{V_0\lambda_V^{-2}} \;=\; \lambda_V^4 .
\label{eq:lambda4}
\end{equation}
\end{theorem}

\begin{proposition}[ROB: Robustness of \eqref{eq:lambda4}]
\label{prop:robust}
The ratio depends only on the saturation limits of $\tanh$, not on
$\theta$, and is therefore invariant to the shape of the transition between
regimes.
\end{proposition}

\begin{proposition}[EGE: EGARCH-class equivalence]
\label{prop:egarch}
The recursion \eqref{eq:zrec} is observationally equivalent to an
EGARCH-class specification (exponential link, bounded, asymmetric) in
$\log V$, conditioned on the state level $z_{t-1}$ rather than on a
standardised residual as in Nelson's (1991) canonical EGARCH. Given
$\theta$, $\lambda_V$ is estimable from returns data without observing $Q_t$.
\end{proposition}

\begin{proposition}[SJT: Sojourn time]
\label{prop:sojourn}
For a buffer with drift $\mu>0$ started a distance $d$ below the
requirement, the expected first-passage time to $z=0$ is
$\E[T] = d/\mu$, with the comparative statics in $d$ and $\mu$ immediate.
\end{proposition}

\noindent
Theorem~\ref{thm:lambda4} is the model's one directly estimable
implication. It is Lean-verified and depends on nothing in
Sections~\ref{sec:formulation}--\ref{sec:numerical}.

\section{The continuous-time formulation}
\label{sec:formulation}

\subsection{State, controls, objective}

Let $X_t$ be the capital buffer and $R$ a reference level. Three controls
of three distinct mathematical types act simultaneously:
\begin{itemize}
\item $\pi_t$, a \emph{classical} control of risky exposure, constrained by
  a state-dependent regulatory cap $\pi \le \bar\pi(x)$;
\item $D_t$, a \emph{singular} control of dividends, acting at an upper
  barrier;
\item $I_t$, an \emph{impulse} control of recapitalisation, with cost
  $K + (1+\kappa)\xi$ for an injection of size $\xi$ --- a fixed cost
  \emph{and} a proportional cost.
\end{itemize}
Shortfall below the reference is penalised at rate
\begin{equation}
\Lambda(x) \;=\; (\lambda_S-1)(R-x)^+ , \qquad \lambda_S \ge 1 .
\label{eq:Lambda}
\end{equation}
The state is the assets-to-liabilities ratio $x = A/L \in (1,\infty)$, and
$V$ is the value \emph{per unit of liabilities}: if $J(A,L)$ denotes the
unreduced value, then $J(A,L) = L\,V(x)$. This normalisation is what makes
the problem one-dimensional, and it has two linked consequences for the
equation below. Since $\partial J/\partial L = V - xV'$, the liability
drift contributes $\mu_L L (V - xV')$ to the unreduced generator; dividing
by $L$, the term $-\mu_L x V'$ joins the drift --- which is why
\eqref{eq:drift} carries $-\mu_L x$ --- and the discount term becomes
\begin{equation}
-\rho V + \mu_L V \;=\; -\rho_L V ,
\qquad \rho_L \;:=\; \rho - \mu_L ,
\label{eq:rhoL}
\end{equation}
so the \emph{liability-net} rate $\rho_L$, not the gross rate $\rho$,
discounts the reduced problem. Both halves of the reduction are forced
together: an equation carrying $-\mu_L x$ in its drift must carry
$-\rho_L V$ in its discount term.

Writing $\mathcal{L}^\pi$ for the controlled generator of $x$, the value
function satisfies the quasi-variational inequality
\begin{equation}
\max\Big\{
\sup_{\pi \le \bar\pi(x)} \mathcal{L}^\pi V(x) - \rho_L V(x) - \Lambda(x),
\;\; 1 - V'(x),
\;\; \mathcal{M}V(x) - V(x)
\Big\} \;=\; 0,
\label{eq:qvi}
\end{equation}
with $\mathcal{M}V(x) = \sup_{\xi>0}\{V(x+\xi) - K - (1+\kappa)\xi\}$.
The drift and diffusion are
\begin{align}
\mu(x,\pi) &= x\big((1-\pi)r + \pi\mu_s - \mu_L\big) + \gamma,
\label{eq:drift}\\
\sigma^2(x,\pi) &= \pi^2\sigma^2 x^2
  + 2\pi c\,\sigma\sigma_L\, x(1-x)
  + \sigma_L^2 (1-x)^2 ,
\label{eq:diff}
\end{align}
where $c$ is the correlation between the risky return and the liability
shock, and $\mu_L = \gamma + r_L$ combines liability growth and the rate
paid on liabilities.

The supremum in \eqref{eq:qvi} is taken over the constrained set
$\pi \le \bar\pi(x)$, so the formulation already accommodates both a
binding and a slack cap: where the unconstrained maximiser lies below
$\bar\pi(x)$ the supremum takes it, and where it does not the supremum
takes $\bar\pi(x)$. No separate boundary condition is attached to the
switch between the two.

\subsection{Two structural properties}

\begin{proposition}[NST: Exact nesting of the risk-neutral benchmark]
\label{prop:nesting}
$\Lambda(x) \equiv 0$ identically at $\lambda_S = 1$. The problem therefore
\emph{nests} the risk-neutral benchmark exactly, rather than approaching it
in a limit.
\end{proposition}

\begin{proposition}[RPC: The running payoff is concave]
\label{prop:concave-payoff}
$-\Lambda$ is concave on $\mathbb{R}$ for every $\lambda_S \ge 1$.
\end{proposition}

\noindent
Proposition~\ref{prop:concave-payoff} marks the departure from
reference-dependent models built on S-shaped utility, which are convex
below the reference and therefore require a concavification argument
before any threshold structure can be recovered: \eqref{eq:Lambda} is
kinked-linear rather than S-shaped, so the running payoff stays concave.
The claim should be stated at exactly that strength and no further. It is
a statement about $\Lambda$, and it does \emph{not} deliver a concave $V$:
Remark~\ref{rem:nonconcave} finds $V$ locally convex above the
recapitalisation boundary, caused by the fixed issuance cost $K$ and
present already at $\lambda_S=1$. The difference from the S-shaped
literature is that the non-concavity here is localised and
cost-driven rather than global and preference-driven.

\subsection{Coefficients where the cap binds}
\label{sec:capcoef}

Where the regulatory cap binds, $\pi = \bar\pi(x)$ and the supremum in
\eqref{eq:qvi} is attained at the constraint, so the generator's
coefficients are
\begin{equation}
A(x) = \tfrac12\,\sigma^2\big(x,\bar\pi(x)\big),
\qquad
B(x) = \bar\pi(x)\,x\,(\mu_s-r) + (r-\mu_L)x + \gamma ,
\label{eq:capcoef}
\end{equation}
with $\sigma^2$ as in \eqref{eq:diff}. Write $u(x) = \bar\pi(x)\,x$ for the
capped exposure, $u(x) = \min\{(x-1)/a_1,\,(x-a_2)/a_3\}$, and let
$\bar x$ denote the level at which the two branches cross. On the solvency
regime $u(x) = (x-1)/a_1$ and \eqref{eq:diff} gives
$\sigma^2(x,\bar\pi(x)) = \kappa_1(c)^2 (x-1)^2$ with
\begin{equation}
\kappa_1(c)^2 \;=\; \frac{\sigma^2}{a_1^2}
   - \frac{2c\,\sigma\sigma_L}{a_1} + \sigma_L^2 .
\label{eq:kappa1c}
\end{equation}

\begin{proposition}[CCP: Consistency condition on parameters]
\label{prop:rrL}
$u(1)=0$, so $B(1) = (r-\mu_L) + \gamma = r - r_L$ under
$\mu_L = \gamma + r_L$. A positive drift at the distress boundary
therefore requires $r > r_L$.
\end{proposition}

\noindent
The parameters used throughout satisfy this: $r=0.025$, $\mu_L=0.03$,
$\gamma=0.02$, giving $r_L=0.01$ and $r-r_L=+0.015>0$.

\begin{proposition}[LQB: Lower bound, the liquidation payoff]
\label{prop:lowerbound}
$V(x) \ge x-1$ for every $x \ge 1$; in particular $V(1) \ge 0$.
\end{proposition}
\begin{proof}
The strategy that pays an immediate lump-sum dividend of $x-1$, driving
the state to the bankruptcy point, is admissible and collects exactly
$x-1$. Because the action is instantaneous, the running penalty
$\int e^{-\rho_L t}\Lambda(X_t)\,dt$ accrues nothing: $\Lambda$ enters the
objective only, never the admissible set, so the bound is unaffected by
$\lambda_S$. Since $V$ is a supremum over admissible strategies,
$V(x)\ge x-1$.
\end{proof}

\section{The state variable, and the boundary at \texorpdfstring{$x=1$}{x=1}}
\label{sec:statemap}

Section~\ref{sec:discrete} uses $z = \log(E/K)$: one special point ($z=0$),
no finite lower boundary, equity exhaustion at $z=-\infty$.
Section~\ref{sec:formulation} uses a level ratio of assets to liabilities,
where $x=1$ means equity exactly zero and the diffusion
\eqref{eq:diff} degenerates. These look like different variables, with two
special points rather than one and a finite degenerate boundary rather
than none. They are not.

\begin{proposition}[TSO: The two states are one state]
\label{prop:statemap}
Let the capital requirement be a ratio requirement, $Q = q L$. Then
$E/L = x-1$ and $E/Q = (x-1)/q$, so
\begin{equation}
z \;=\; \log\!\frac{x-1}{q},
\qquad
x \;=\; 1 + q\,e^{z},
\label{eq:statemap}
\end{equation}
a diffeomorphism of $(1,\infty)$ onto $\mathbb{R}$. Under it:
\begin{enumerate}
\item[(i)] the special points coincide --- $x=1 \leftrightarrow
z=-\infty$, $x=1+q \leftrightarrow z=0$, $x\to\infty \leftrightarrow
z\to\infty$;
\item[(ii)] on the solvency regime the volatility of $z$ is the
\emph{constant} $\kappa_1(c)$ of \eqref{eq:kappa1c}, for every
correlation $c$;
\item[(iii)] the drift of $z$ is $B(x)/(x-1) - \kappa_1^2/2$, and
$B(1) = r-r_L$ by Proposition~\ref{prop:rrL}, so it tends to $+\infty$ as
$z\to-\infty$ precisely when $r>r_L$.
\end{enumerate}
\end{proposition}

\begin{corollary}[CDA: The degeneracy is a coordinate artefact, and the
boundary is unreachable]
\label{cor:boundary}
By (ii) the vanishing of the diffusion at $x=1$ is exactly cancelled by
the blow-up of $\dd z/\dd x$: in $z$ the diffusion is non-degenerate
everywhere, so the degeneracy is a property of the imported coordinate,
not of the model. By (iii), in $z$-coordinates the process has constant
volatility and a restoring drift growing without bound as
$z \to -\infty$ whenever $r>r_L$; equity is therefore never exhausted in
finite time, and $x=1$ is unreachable. Unreachability is
coordinate-invariant --- $x=1$ is unreachable if and only if $z=-\infty$
is --- while its presentation as a degenerate-diffusion problem is not.
\end{corollary}

\noindent
Two limitations, stated rather than buried. The map requires reading
Section~\ref{sec:discrete}'s exogenous $K$ as a \emph{ratio} requirement
set by regulation rather than a deterministic constant; with $K$ constant
and $L$ stochastic no deterministic map exists. And (ii) holds on the
solvency regime only --- above $\bar x$ the cap switches to
$(x-a_2)/a_3$ and the $z$-volatility is not constant. That range covers
where the boundary question lives, since the recapitalisation boundaries
of Section~\ref{sec:numerical} all lie below $\bar x$, but it is a
restriction, not a global result.

\medskip\noindent
One consequence is worth recording because it is what makes $\lambda_V$
and $\lambda_S$ comparable at all: at the numerical baseline $R=1.15$ the
requirement level is $q = R-1 = 0.15$, and $z=0$ maps to $x=1.15=R$
exactly. The two sections' reference points are the same point.

\begin{proposition}[SCG: The surplus limit is cap geometry too]
\label{prop:satlimits}
On the liquidity regime ($x>\bar x$, $u(x) = (x-a_2)/a_3$), the
$z$-volatility tends, as $x\to\infty$, to the constant
\begin{equation}
\kappa_3(c)^2 \;=\; \frac{\sigma^2}{a_3^2} - \frac{2c\,\sigma\sigma_L}{a_3}
+ \sigma_L^2 ,
\label{eq:kappa3c}
\end{equation}
the same functional form as $\kappa_1(c)^2$ of \eqref{eq:kappa1c} with
$a_1$ replaced by $a_3$. Consequently
\begin{equation}
\lambda_V^4 \;=\; \frac{\kappa_1(c)^2}{\kappa_3(c)^2},
\label{eq:lamV-geom}
\end{equation}
a function of $(\sigma,\sigma_L,a_1,a_3,c)$ only.
\end{proposition}
\begin{proof}
$\bar\pi(x) \to 1/a_3$ as $x\to\infty$; substituting into
\eqref{eq:diff} and taking the limit gives \eqref{eq:kappa3c} directly,
by the same computation as Proposition~\ref{prop:statemap}(ii) with
$a_1\to a_3$. Equation \eqref{eq:lamV-geom} is Theorem~\ref{thm:lambda4}'s
definition evaluated at the two saturation limits.
\end{proof}

\noindent
The content of Proposition~\ref{prop:satlimits} is what
\eqref{eq:lamV-geom} does \emph{not} contain: no $\lambda_S$, $K$,
$\kappa$, $R$, $\rho$, $r$, $\mu_s$, $\mu_L$ or $\gamma$. Evaluated at the
saturation limits, $\lambda_V$ is a function of cap geometry and
correlation alone. At the numerical baseline $\kappa_1=1.7720$ and
$\kappa_3=0.2623$, giving $\lambda_V^4 = 45.633$ and $\lambda_V=2.599$
--- a constant, carrying no information about $\lambda_S$. Two
qualifications bound the reach of this. The limits are approached only as
$x\to1^+$ and $x\to\infty$, and the process's actual range under the
solved policy sits far from both: on a $y_{\max}=2.5$ grid the
surplus-side approach to $\kappa_3$ is still an order of magnitude short
at the upper edge. And \eqref{eq:lamV-geom} is the \emph{definitional}
evaluation of $\lambda_V$; it does not preclude a relation between an
\emph{estimated} $\lambda_V$ and $\lambda_S$ (Section~\ref{sec:open}).

\section{Threshold comparative statics}
\label{sec:compstat-thm}

Both boundaries respond to $\lambda_S$, and both responses reduce to
properties of a single object: the expected discounted shortfall incurred
under the optimal policy,
\begin{equation}
W(x) \;:=\; \mathbb{E}_x\!\left[\int_0^\infty e^{-\rho_L t}\,(R-X_t)^+\,dt\right],
\label{eq:Wdef}
\end{equation}
the expectation taken under the policy optimal at the prevailing
$\lambda_S$. Because $\lambda_S$ enters the objective only through
$\Lambda(x) = (\lambda_S-1)(R-x)^+$, and enters it linearly, $W$ is
exactly the sensitivity of value to the preference parameter.

\begin{lemma}[REG: Regularity in the parameters]
\label{lem:reg}
Let $\theta$ denote either $\lambda_S$ or $K$, the other being held fixed.
Then for every $x$:
\begin{enumerate}
\item[(i)] $\theta \mapsto V(x;\theta)$ is convex and non-increasing;
\item[(ii)] it is therefore locally Lipschitz, admits one-sided
  derivatives at every $\theta$, and is differentiable at all but
  countably many $\theta$;
\item[(iii)] at every $\theta$ the one-sided derivatives are
  $\partial^{+}_\theta V = -\min_{a \in \mathcal{A}^*(\theta)} T(x,a)$ and
  $\partial^{-}_\theta V = -\max_{a \in \mathcal{A}^*(\theta)} T(x,a)$,
  where $\mathcal{A}^*(\theta)$ is the set of optimal policies and
  $T = S$ for $\theta=\lambda_S$, $T = N$ for $\theta = K$; in particular
  $V$ is differentiable at $\theta$ if and only if every optimal policy
  yields the same $T$;
\item[(iv)] $W(x;\cdot)$ is non-increasing in $\lambda_S$ and
  $N(x;\cdot)$ is non-increasing in $K$.
\end{enumerate}
\end{lemma}

\begin{proof}
For fixed $a$, $\theta \mapsto J(x,a;\theta)$ is affine with slope
$-T(x,a) \le 0$: for $\lambda_S$ this is the decomposition
$J = A - (\lambda_S-1)S$ used below, and for $K$ it is $J = B - KN$. Both
are exact because the parameter multiplies a single term of the objective
and neither restricts the admissible set. $V(x;\cdot)$ is the pointwise
supremum of this family, and a supremum of affine functions is convex; the
slopes are non-positive, so the supremum is non-increasing. This gives
(i), and (ii) is standard for a finite convex function on an interval.
Part (iii) is the subdifferential formula for a supremum of affine
functions, $\partial_\theta V = -\operatorname{co}\{T(x,a) : a \in
\mathcal{A}^*(\theta)\}$, whose endpoints are the stated one-sided
derivatives; differentiability is the statement that this interval is a
single point. Part (iv) is (i) restated: convexity means
$\partial_\theta V$ is non-decreasing, and $\partial_\theta V = -T$ at
points of differentiability, so $T$ is non-increasing.
\end{proof}

\noindent
Part (iv) is a testable by-product rather than a technical step, and the
solve bears it out: across the $\lambda_S$ sweep the second divided
differences of $V$ in $\lambda_S$ are strictly positive at every grid
point, and $W$ is strictly decreasing in $\lambda_S$ at every grid point.
Neither was imposed on the solver.

\begin{lemma}[ENV: The envelope identity]
\label{lem:envelope}
At any $\lambda_S$ at which $V(x;\cdot)$ is differentiable --- by
Lemma~\ref{lem:reg}(ii), all but countably many --- and with the optimal
policy measurable in $(x,\lambda_S)$,
\begin{equation}
\frac{\partial V}{\partial \lambda_S}(x) \;=\; -\,W(x),
\label{eq:envelope}
\end{equation}
and $W$ satisfies, on the inaction region,
\begin{equation}
\tfrac12\sigma^2(x,u(x))W''(x) + \mu(x,u(x))W'(x) - \rho_L W(x)
  + (R-x)^+ \;=\; 0 .
\label{eq:Wode}
\end{equation}
\end{lemma}

\begin{proof}
For an admissible policy $a$ write the payoff as
$J(x,a;\lambda_S) = A(x,a) - (\lambda_S-1)S(x,a)$, where
$S(x,a) = \mathbb{E}_x[\int_0^\infty e^{-\rho_L t}(R-X_t)^+dt]$ is the
discounted shortfall under $a$ and $A$ collects the terms not involving
$\lambda_S$ --- dividends received and impulse costs paid. The
decomposition is exact because $\Lambda$ is linear in $\lambda_S$ and
enters the objective only, never the admissible set (as in the proof of
Proposition~\ref{prop:lowerbound}). Thus $\lambda_S \mapsto J(x,a;\lambda_S)$
is affine with slope $-S(x,a)$ for every fixed $a$, and $V$ is their
pointwise supremum. Under the stated hypotheses Danskin's theorem applies
and the derivative of the supremum is the partial derivative evaluated at
the maximiser, giving \eqref{eq:envelope} with $W(x) = S(x,a^*)$.

For \eqref{eq:Wode}, differentiate the interior equation
$\sup_\pi \mathcal{L}^\pi V - \rho_L V - \Lambda = 0$ with respect to
$\lambda_S$. The supremum is attained at $u(x)$, so by the same envelope
argument the derivative of the supremum term is
$\mathcal{L}^{u(x)}(\partial V/\partial\lambda_S)$; and
$\partial\Lambda/\partial\lambda_S = (R-x)^+$. Substituting
$\partial V/\partial\lambda_S = -W$ and negating gives \eqref{eq:Wode}.
\end{proof}

\noindent
The two smooth-fit conditions then convert \eqref{eq:Wode} into boundary
data for $W$. At the payout barrier, $V'(y^*(\lambda_S);\lambda_S)=1$
holds identically along the boundary, so differentiating and using
$V''(y^*)=0$ gives $\partial V'/\partial\lambda_S(y^*)=0$, that is
\begin{equation}
W'(y^*) \;=\; 0 .
\label{eq:Wneumann}
\end{equation}
At the recapitalisation trigger, differentiating the value-matching
condition $V(x_L) = \mathcal{M}V(x_L)$ along the boundary, with the
first-order terms cancelling by the impulse smooth fit
$V'(x_L) = (\mathcal{M}V)'(x_L) = 1+\kappa$ and the $y_{\mathrm{post}}$
term vanishing by the envelope property of the inner supremum, gives the
nonlocal condition
\begin{equation}
W(x_L) \;=\; W(y_{\mathrm{post}}) .
\label{eq:Wnonlocal}
\end{equation}
The second-order smooth fit $V''(y^*)=0$ is what makes the barrier
tractable rather than what obstructs it: it is precisely the condition
delivering the Neumann datum \eqref{eq:Wneumann}.

\begin{lemma}[BDR: Differentiability of the boundaries]
\label{lem:bdr}
Let $\theta$ be either parameter. Suppose $V$ is $C^3$ in $x$ in a
neighbourhood of $y^*$ and $C^2$ in $x$ near $x_L$ and
$y_{\mathrm{post}}$, with those $x$-derivatives jointly continuous in
$(x,\theta)$ and continuously differentiable in $\theta$. Then, on the set
where the corresponding non-degeneracy condition holds,
\begin{enumerate}
\item[(i)] $y^*(\theta)$ is continuously differentiable wherever
  $\rho_L \ne \mu'(y^*)$;
\item[(ii)] $x_L(\theta)$ is continuously differentiable wherever
  $V''(x_L) \ne 0$;
\item[(iii)] $y_{\mathrm{post}}(\theta)$ is continuously differentiable
  wherever $V''(y_{\mathrm{post}}) \ne 0$.
\end{enumerate}
\end{lemma}

\begin{proof}
Each boundary is the zero set of a scalar function to which the implicit
function theorem applies; in each case the required non-degeneracy is the
$x$-derivative of that function.

For $y^*$, the defining relation is not $V'(y)-1=0$ --- that equation is
degenerate here, since $V''(y^*)=0$ --- but the second-order smooth fit
itself. Set $F(y,\theta) := V''(y;\theta)$, so $F(y^*(\theta),\theta)=0$.
Then $\partial F/\partial y = V'''(y^*)$, and evaluating the differentiated
interior equation at $y^*$, where $V''=0$, $V'=1$ and $\Lambda'=0$, gives
$\tfrac12\sigma^2(y^*)V'''(y^*) = \rho_L - \mu'(y^*)$. Since
$\sigma^2(y^*)>0$, $\partial F/\partial y \ne 0$ exactly when
$\rho_L \ne \mu'(y^*)$, which is (i).

For $x_L$, set $G(x,\theta) := V'(x;\theta) - (1+\kappa)$, so
$G(x_L(\theta),\theta) = 0$ by the impulse smooth fit, and
$\partial G/\partial x = V''(x_L)$, giving (ii). For
$y_{\mathrm{post}}$ the same function serves, with the first-order
condition $V'(y_{\mathrm{post}}) = 1+\kappa$ in place of the smooth fit
and $\partial G/\partial x = V''(y_{\mathrm{post}})$, giving (iii).
\end{proof}

\noindent
The three non-degeneracy conditions in Lemma~\ref{lem:bdr} are not new
hypotheses: they are the same conditions --- $\rho_L > \mu'(y^*)$,
$V''(x_L)>0$, $V''(y_{\mathrm{post}})<0$ --- that sign the derivatives in
Propositions~\ref{prop:tcs} and~\ref{prop:kcs}. Each therefore does double
duty, delivering both the existence of the derivative and its sign. What
remains genuinely assumed is the joint regularity of $V$ in $(x,\theta)$
stated at the head of the lemma, which is a regularity question for the
quasi-variational inequality itself and is not settled here.

\begin{proposition}[TCS: Threshold comparative statics]
\label{prop:tcs}
Assume Lemmas~\ref{lem:reg}--\ref{lem:bdr}, work at a $\lambda_S$ of
differentiability, and take $V''(x_L)>0$ from
Proposition~\ref{prop:soc}(ii). Then
\begin{equation}
\frac{dy^*}{d\lambda_S} \;=\; \frac{\rho_L\,W(y^*)}{\rho_L - \mu'(y^*)},
\qquad
\frac{dx_L}{d\lambda_S} \;=\; \frac{W'(x_L)}{V''(x_L)} .
\label{eq:tcs}
\end{equation}
Consequently:
\begin{enumerate}
\item[(i)] $dy^*/d\lambda_S > 0$ whenever $W(y^*)>0$ and
  $\rho_L > \mu'(y^*)$, the latter being equivalent to
  $\rho > r + \pi(\mu_s - r)$ at $y^*$;
\item[(ii)] $dx_L/d\lambda_S > 0$ whenever
  $\rho_L\,W(y_{\mathrm{post}}) < R - x_L$.
\end{enumerate}
\end{proposition}

\begin{proof}
\emph{Barrier.} Differentiating $V'(y^*(\lambda_S);\lambda_S)=1$ gives
$V''(y^*)\,dy^*/d\lambda_S + \partial V'/\partial\lambda_S(y^*) = 0$,
which is $0=0$ by $V''(y^*)=0$ and \eqref{eq:Wneumann}: the first-order
condition is degenerate and carries no information. Differentiate instead
the second-order condition $V''(y^*(\lambda_S);\lambda_S)=0$, giving
$V'''(y^*)\,dy^*/d\lambda_S = W''(y^*)$, since
$\partial V''/\partial\lambda_S = -W''$. Evaluate the two third-order
quantities. In \eqref{eq:Wode} at $y^*$, using $W'(y^*)=0$ and
$(R-y^*)^+=0$ for $y^*>R$, one gets
$\tfrac12\sigma^2(y^*)W''(y^*) = \rho_L W(y^*)$. Differentiating the
interior equation and evaluating at $y^*$, where $V''=0$, $V'=1$ and
$\Lambda'=0$, gives
$\tfrac12\sigma^2(y^*)V'''(y^*) = \rho_L - \mu'(y^*)$. Taking the ratio,
the common factor $\tfrac12\sigma^2(y^*)$ cancels and the first identity
in \eqref{eq:tcs} follows. For (i), $W\ge0$ always and $W(y^*)>0$ whenever
the buffer reaches deficit with positive probability; and since
$\mu'(y^*) = \mu_{\mathrm{opt}}(y^*) - \mu_L$ with
$\mu_{\mathrm{opt}} = r + \pi(\mu_s-r)$, the condition
$\rho_L > \mu'(y^*)$ is $\rho > \mu_{\mathrm{opt}}(y^*)$.

\emph{Trigger.} Differentiating the impulse smooth fit
$V'(x_L(\lambda_S);\lambda_S) = 1+\kappa$ gives
$V''(x_L)\,dx_L/d\lambda_S + \partial V'/\partial\lambda_S(x_L) = 0$,
and $\partial V'/\partial\lambda_S = -W'$, yielding the second identity in
\eqref{eq:tcs}. For the sign, let $\tau$ be the first hitting time of
$x_L$ from $x>x_L$ and decompose \eqref{eq:Wdef} at $\tau$, using that the
policy injects at $x_L$ and restarts from $y_{\mathrm{post}}$:
\[
W(x) - W(y_{\mathrm{post}})
= \mathbb{E}_x\!\left[\int_0^\tau e^{-\rho_L t}(R-X_t)^+dt\right]
  - \big(1-\mathbb{E}_x[e^{-\rho_L\tau}]\big)W(y_{\mathrm{post}}).
\]
As $x \downarrow x_L$ we have $\tau \to 0$, and to first order in
$\mathbb{E}_x[\tau]$ the right-hand side is
$\mathbb{E}_x[\tau]\,\big[(R-x_L) - \rho_L W(y_{\mathrm{post}})\big] + o(\mathbb{E}_x[\tau])$,
using $\mathbb{E}_x[\int_0^\tau e^{-\rho_L t}(R-X_t)^+dt] = (R-x_L)\mathbb{E}_x[\tau] + o(\mathbb{E}_x[\tau])$
and $1-\mathbb{E}_x[e^{-\rho_L\tau}] = \rho_L\mathbb{E}_x[\tau] + o(\mathbb{E}_x[\tau])$.
Since $\mathbb{E}_x[\tau]$ is strictly increasing in $x$ near $x_L$ and
vanishes at $x_L$, the sign of $W'(x_L)$ is that of
$(R-x_L) - \rho_L W(y_{\mathrm{post}})$. With $V''(x_L)>0$ this gives (ii).
\end{proof}

\begin{remark}[TCSN: What the proposition rests on]
\label{rem:tcsn}
The differentiability of $V$ in $\lambda_S$ and of the two boundaries is
established, not assumed: the first by Lemma~\ref{lem:reg} (the
value is convex in each parameter, hence differentiable off a countable
set), the second by Lemma~\ref{lem:bdr} (the implicit function theorem,
whose non-degeneracy conditions are the sign conditions below). The
hypothesis $V''(x_L)>0$ is Proposition~\ref{prop:soc}(ii), which supplies
it from the variational inequality up to the non-degeneracy $\Psi(x_L)>0$;
it is no longer the numerical finding of Remark~\ref{rem:nonconcave}. What
the solve adds is not confirmation of the identities --- $W$ is measured
there as $-\partial V/\partial\lambda_S$, so \eqref{eq:envelope} is
definitional in the numerics and the barrier identity is a rearrangement
of the barrier condition itself --- but a check of the two sign
conditions. At the stated parameters $\rho_L - \mu'(y^*) = 0.045 > 0$, and
condition~(i) is implied by the parameter-validity requirement
$\rho > r + \mu_r^+/\bar a$ already imposed on the solve, since the cap
bounds $\pi(\mu_s-r) \le \mu_r^+/\bar a$. Condition~(ii) holds across the
sweep with $R-x_L$ exceeding $\rho_L W(y_{\mathrm{post}})$ by a factor
between $11$ and $28$, the margin widening in $\lambda_S$. Both are
therefore slack rather than marginal at the calibration, though neither is
guaranteed in general: (ii) fails when recapitalisation barely reduces
continuation shortfall, and in that regime the trigger need not rise.
\end{remark}

\subsection{The fixed cost, and joint identification}
\label{sec:kcs}

The same construction applies to the fixed issuance cost. Like $\lambda_S$,
the parameter $K$ enters the objective linearly and never enters the
admissible set, so the argument of Lemma~\ref{lem:envelope} transfers
verbatim with one substitution: the running shortfall is replaced by the
discounted \emph{count} of injections,
\begin{equation}
N(x) \;:=\; \mathbb{E}_x\!\left[\sum_{i} e^{-\rho_L \tau_i}\right],
\label{eq:Ndef}
\end{equation}
the sum over injection times $\tau_i$ under the optimal policy.

\begin{lemma}[ENVK: The envelope identity for the fixed cost]
\label{lem:envelopeK}
At any $K$ at which $V(x;\cdot)$ is differentiable --- by
Lemma~\ref{lem:reg}(ii), all but countably many ---
\begin{equation}
\frac{\partial V}{\partial K}(x) \;=\; -\,N(x),
\label{eq:envelopeK}
\end{equation}
and on the inaction region $N$ satisfies the \emph{homogeneous} equation
\begin{equation}
\tfrac12\sigma^2(x,u(x))N''(x) + \mu(x,u(x))N'(x) - \rho_L N(x) \;=\; 0 ,
\label{eq:Node}
\end{equation}
with boundary data $N'(y^*)=0$ and $N(x_L) = N(y_{\mathrm{post}}) + 1$.
\end{lemma}

\begin{proof}
Write $J(x,a;K) = B(x,a) - K\,N(x,a)$, where $N(x,a)$ is the discounted
injection count under $a$ and $B$ collects the terms not involving $K$.
The decomposition is exact because $K$ multiplies the number of impulses
and nothing else, and it does not restrict the admissible set; Danskin's
theorem then gives \eqref{eq:envelopeK} as in Lemma~\ref{lem:envelope}.
Equation \eqref{eq:Node} is homogeneous rather than inhomogeneous because
injections accrue only at boundary hits: no injection occurs while the
state is interior, so there is no running term. The Neumann datum at $y^*$
follows from differentiating $V'(y^*(K);K)=1$ and using $V''(y^*)=0$,
exactly as \eqref{eq:Wneumann} did. At the trigger, differentiate
$V(x_L) = V(y_{\mathrm{post}}) - K - (1+\kappa)(y_{\mathrm{post}}-x_L)$
along the boundary; the terms in $y_{\mathrm{post}}'$ cancel by the
first-order condition $V'(y_{\mathrm{post}})=1+\kappa$, the terms in
$x_L'$ cancel by the smooth fit $V'(x_L)=1+\kappa$, and the explicit
$-\partial K/\partial K = -1$ survives, giving
$\partial V/\partial K(x_L) = \partial V/\partial K(y_{\mathrm{post}}) - 1$,
which is the stated condition. It also follows directly from
\eqref{eq:Ndef}: from $x_L$ the policy injects at once and then faces the
continuation count from $y_{\mathrm{post}}$.
\end{proof}

\begin{proposition}[KCS: Fixed-cost comparative statics]
\label{prop:kcs}
Assume Lemma~\ref{lem:envelopeK} and Lemmas~\ref{lem:reg}--\ref{lem:bdr},
work at a $K$ of differentiability, and take $V''(x_L)>0$ and
$V''(y_{\mathrm{post}})<0$ from Proposition~\ref{prop:soc}. Then
\begin{equation}
\frac{dy^*}{dK} = \frac{\rho_L\,N(y^*)}{\rho_L - \mu'(y^*)} > 0,
\qquad
\frac{dx_L}{dK} = \frac{N'(x_L)}{V''(x_L)} < 0,
\qquad
\frac{dy_{\mathrm{post}}}{dK} = \frac{N'(y_{\mathrm{post}})}{V''(y_{\mathrm{post}})} > 0,
\label{eq:kcs}
\end{equation}
the first under $\rho_L > \mu'(y^*)$. In particular the inaction region
widens at both ends, and the trigger-to-target gap
$y_{\mathrm{post}} - x_L$ is strictly increasing in $K$.
\end{proposition}

\begin{proof}
For $x$ above the trigger, conditioning on the first hitting time $\tau$
of $x_L$ gives
$N(x) = \mathbb{E}_x[e^{-\rho_L\tau}]\big(1 + N(y_{\mathrm{post}})\big)$,
and $x \mapsto \mathbb{E}_x[e^{-\rho_L\tau}]$ is strictly decreasing, so
$N' < 0$ throughout the inaction region and $N > 0$. The barrier identity
is obtained exactly as in Proposition~\ref{prop:tcs}: the first-order
condition is degenerate, and differentiating $V''(y^*(K);K)=0$ gives
$V'''(y^*)\,dy^*/dK = N''(y^*)$; evaluating \eqref{eq:Node} at $y^*$ with
$N'(y^*)=0$ gives $\tfrac12\sigma^2 N''(y^*) = \rho_L N(y^*)$, and
$\tfrac12\sigma^2 V'''(y^*) = \rho_L - \mu'(y^*)$ as before. The trigger
identity follows from differentiating the smooth fit
$V'(x_L(K);K) = 1+\kappa$, and is negative since $N'(x_L)<0$ and
$V''(x_L)>0$. The target identity follows from differentiating the
first-order condition $V'(y_{\mathrm{post}}(K);K) = 1+\kappa$, and is
positive since $N'(y_{\mathrm{post}})<0$ and $V''(y_{\mathrm{post}})<0$
in the concave interior. The gap statement is the difference of the last
two.
\end{proof}

\begin{corollary}[IDN: Joint local identification]
\label{cor:idn}
Under Propositions~\ref{prop:tcs} and~\ref{prop:kcs} and their sign
conditions, the Jacobian of $(y^*,x_L)$ with respect to
$(\lambda_S,K)$ is nonsingular:
\begin{equation}
\det \frac{\partial(y^*,x_L)}{\partial(\lambda_S,K)}
= \frac{\rho_L}{\big(\rho_L-\mu'(y^*)\big)V''(x_L)}
  \Big[\,W(y^*)N'(x_L) \;-\; N(y^*)W'(x_L)\,\Big] \;<\; 0 .
\label{eq:jac}
\end{equation}
Hence $(\lambda_S,K)$ is locally identified from the pair of thresholds.
\end{corollary}

\begin{proof}
The matrix entries are those of \eqref{eq:tcs} and \eqref{eq:kcs}; the
common factors $\rho_L/(\rho_L-\mu'(y^*))$ and $1/V''(x_L)$ factor out of
the first row and second row respectively. In the bracket,
$W(y^*)>0$ and $N'(x_L)<0$ make the first product negative, while
$N(y^*)>0$ and $W'(x_L)>0$ make the second positive, so the bracket is
strictly negative; the prefactor is positive. The determinant is therefore
strictly negative, in particular nonzero, and the inverse function theorem
applies.
\end{proof}

\begin{remark}[KCSN: Reading the identification result]
\label{rem:kcsn}
The two parameters separate for a structural reason rather than a
numerical one: $\lambda_S$ and $K$ both raise the payout barrier, but they
move the recapitalisation trigger in \emph{opposite} directions --- a
steeper shortfall penalty makes the firm intervene earlier, a larger fixed
cost makes it wait longer. That opposition is what signs \eqref{eq:jac}
away from zero, and it is why the determinant does not depend on comparing
magnitudes. The claim is local identification from the thresholds, not a
statement that any estimator attains it.

Two cautions, and one piece of support. The result inherits every
hypothesis of the two propositions, including $V''(x_L)>0$, now
Proposition~\ref{prop:soc}(ii) rather than a numerical finding; and local
identification from the thresholds is not a claim that any particular
estimator attains it. The support is that the converged $K$-sweep of
Remark~\ref{rem:Ksens}, run at $\lambda_S=1$ across
$K\in\{0,\dots,0.05\}$ with every run converging, already exhibits both
predicted directions. The continuation band $y^*-x_L$ widens
monotonically from $0.159$ to $0.685$, which is
$dy^*/dK - dx_L/dK > 0$; and at $K=0.05$ the recapitalisation region is
one cell wide with the trigger collapsed onto the lower boundary, which
is $dx_L/dK<0$ directly. Combining the two --- a band of $0.685$ with the
trigger at the boundary places $y^*$ near $1.69$, against $1.5715$ at
$K=0.02$ --- indicates $dy^*/dK>0$ as well, though that last step chains
two separately reported quantities rather than reading a recorded $y^*$,
and a sweep recording both thresholds directly would settle it without
the inference.
\end{remark}

\section{Second-order conditions at the impulse boundaries}
\label{sec:secondorder}

The hypotheses $V''(x_L)>0$ and $V''(y_{\mathrm{post}})<0$ have so far been
carried as numerical findings. Both are in fact second-order conditions,
and both follow from the structure of the impulse operator once one
observation is made: $\mathcal{M}v$ is \emph{affine}. Writing the injection
in terms of the target $y = x+\xi$,
\begin{equation}
\mathcal{M}v(x) \;=\; \sup_{y>x}\big\{V(y)-(1+\kappa)(y-x)\big\} - K
\;=\; (1+\kappa)x - K + \sup_{y}\big\{V(y)-(1+\kappa)y\big\},
\label{eq:Maffine}
\end{equation}
the supremum on the right being independent of $x$. Two consequences
follow at once. The injection target does not depend on where the
injection is made from, so $y_{\mathrm{post}}$ is a single number rather
than a function of $x$; and $g := \mathcal{M}v - V$ satisfies
$g' = (1+\kappa)-V'$ and $g'' = -V''$.

Write $\varphi(y) := V(y)-(1+\kappa)y$, so that $y_{\mathrm{post}}$
maximises $\varphi$ and $\varphi' = V'-(1+\kappa)$. It is convenient to
record the sign test the interior equation supplies. At any point where
$V'=1+\kappa$, substituting into
$\tfrac12\sigma^2V''+\mu V'-\rho_LV-\Lambda=0$ gives
\begin{equation}
\tfrac12\sigma^2(x,u(x))\,V''(x) \;=\; \Psi(x),
\qquad
\Psi(x) \;:=\; \rho_L V(x) + \Lambda(x) - (1+\kappa)\mu(x,u(x)),
\label{eq:Psi}
\end{equation}
so at such a point $V''$ and $\Psi$ carry the same sign.

\begin{proposition}[SOC: Second-order conditions]
\label{prop:soc}
Assume $V'>1+\kappa$ on $(x_L,y_{\mathrm{post}})$ and $V'<1+\kappa$ on
$(y_{\mathrm{post}},y^*)$, and, for the second half of (ii) only, the
smooth fit $V'(x_L)=1+\kappa$ --- which is
Proposition~\ref{prop:smf}, proved independently below.
Then
\begin{enumerate}
\item[(i)] $V''(y_{\mathrm{post}}) < 0$, strictly;
\item[(ii)] $V''(x_L^{+}) \ge 0$, with equality if and only if the
  interior equation holds with equality at $x_L$ from the intervention
  side; equivalently, $V''(x_L^{+})>0$ if and only if $\Psi(x_L)>0$.
\end{enumerate}
\end{proposition}

\begin{proof}
(i) By hypothesis $\varphi' > 0$ on $(x_L,y_{\mathrm{post}})$ and
$\varphi'<0$ on $(y_{\mathrm{post}},y^*)$, so $\varphi'$ changes sign at
$y_{\mathrm{post}}$ and its zero there has odd order. If
$\varphi''(y_{\mathrm{post}})=0$ then, $\varphi$ being $C^3$ near
$y_{\mathrm{post}}$ by Lemma~\ref{lem:itr}, $\varphi'$ would vanish to
order at least two, and a zero of even order admits no sign change --- a
contradiction unless $\varphi'$ vanishes identically nearby, which the
hypothesis excludes. Hence $\varphi''(y_{\mathrm{post}}) \ne 0$; being a
maximum it satisfies $\varphi'' \le 0$, so
$V''(y_{\mathrm{post}}) = \varphi''(y_{\mathrm{post}}) < 0$.

(ii) On the intervention region $V = \mathcal{M}V$, which by
\eqref{eq:Maffine} is affine with slope $1+\kappa$; hence $V'=1+\kappa$
and $V''=0$ there. The QVI requires the first branch of \eqref{eq:qvi} to
be non-positive wherever the impulse branch is active, so for $x \le x_L$
\[
\tfrac12\sigma^2\cdot 0 + \mu(x)(1+\kappa) - \rho_L V(x) - \Lambda(x)
\;\le\; 0,
\]
which is exactly $\Psi(x) \ge 0$ on the intervention region, and by
continuity $\Psi(x_L)\ge0$. This much uses only the variational
inequality on the intervention side; it is what
Proposition~\ref{prop:smf} consumes, and it is why that proposition does
not depend on the remainder of this one. Since $V'(x_L)=1+\kappa$ by
Proposition~\ref{prop:smf},
\eqref{eq:Psi} applies at $x_L$ and gives
$\tfrac12\sigma^2 V''(x_L^+) = \Psi(x_L) \ge 0$, so $V''(x_L^+)\ge0$ with
equality precisely when $\Psi(x_L)=0$, that is when the displayed
inequality binds. The same conclusion follows without \eqref{eq:Psi} from
$g\le0$, $g(x_L)=0$ and $g'(x_L)=0$, which force $g''(x_L^+)\le0$.
\end{proof}

\begin{remark}[SOCN: Strictness at the trigger]
\label{rem:socn}
Part (i) is unconditional given the hypotheses: nothing numerical enters.
Part (ii) delivers the weak inequality unconditionally and locates
strictness in a single interpretable condition --- that waiting is
\emph{strictly} suboptimal just below the trigger, rather than marginally
so. That is the natural non-degeneracy for an intervention region with
non-empty interior, and it fails only in the knife-edge case where the
firm is indifferent between acting and waiting on a whole neighbourhood.
At the numerical baseline $\Psi$ is bounded away from zero throughout the
intervention region, with $\Psi(x_L)$ equal to $0.012$, $0.023$ and
$0.035$ at $\lambda_S = 1$, $1.25$ and $2$ respectively, rising in
$\lambda_S$; the corresponding $V''(x_L^+)$ are $6.4$, $13.9$ and $18.7$,
of the same sign in every case as \eqref{eq:Psi} requires.

The consequence for the results above is that
Remark~\ref{rem:nonconcave} is no longer load-bearing.
Propositions~\ref{prop:tcs} and~\ref{prop:kcs},
Corollary~\ref{cor:idn}, Lemma~\ref{lem:bdr} and
Proposition~\ref{prop:ver} took $V''(x_L)>0$ and
$V''(y_{\mathrm{post}})<0$ as hypotheses supported by a numerical
finding; they now rest on Proposition~\ref{prop:soc}, which supplies the
first up to a non-degeneracy condition and the second outright. What
those results still inherit is the smooth fit at $x_L$ --- no longer an
assumption either, but Proposition~\ref{prop:smf} below.
\end{remark}

\begin{proposition}[SMF: Smooth fit at the trigger]
\label{prop:smf}
Let $V$ be the value function and suppose $\sigma^2(x_L,\pi)>0$ for the
optimal $\pi$ at $x_L$, which holds since $x_L>1$
(Lemma~\ref{lem:itr}). Then $V$ is differentiable at $x_L$, with
\[
V'(x_L^-) \;=\; V'(x_L^+) \;=\; 1+\kappa .
\]
\end{proposition}

\begin{proof}
Write $J := V'(x_L^+)-V'(x_L^-)$, both one-sided derivatives existing by
convexity of $V-\mathcal{M}V$ near $x_L$ and interior regularity above it.

\emph{Step 1: $J \ge 0$, and the left derivative is exactly $1+\kappa$.}
On the intervention region $V=\mathcal{M}V$, which by \eqref{eq:Maffine} is
affine with slope $1+\kappa$; hence $V'(x_L^-)=1+\kappa$. The obstacle
$g:=V-\mathcal{M}V$ satisfies $g\ge0$ with $g(x_L)=0$, so $x_L$ minimises
$g$ from the right and $g'(x_L^+)\ge0$, that is $V'(x_L^+)\ge1+\kappa$.
Any kink at $x_L$ is therefore convex. Note this uses only the variational
inequality, not smooth fit, so no circularity arises.

\emph{Step 2: $J>0$ contradicts the supersolution property.}
$V$ is a viscosity solution of \eqref{eq:qvi}; in particular a
supersolution, so for every $\varphi\in C^2$ touching $V$ from below at
$x_L$,
\[
\min\Big\{-\big(\mathcal{L}^\pi\varphi-\rho_LV-\Lambda\big),\;
\varphi'-1,\; V-\mathcal{M}V\Big\}(x_L) \;\ge\; 0 .
\]
Since $V(x_L)=\mathcal{M}V(x_L)$ the third entry vanishes, so the first
must be non-negative. Suppose $J>0$ and write
$\varphi(x)=V(x_L)+p(x-x_L)+\tfrac12 q(x-x_L)^2+o\big((x-x_L)^2\big)$.
For $\varphi\le V$ near $x_L$ it suffices that
$p\in\big(1+\kappa,\,1+\kappa+J\big)$: on each side $V-\varphi$ then has a
strictly signed linear term, which dominates the quadratic for small
$|x-x_L|$ \emph{whatever the value of $q$}. Thus $q$ may be taken
arbitrarily large. But
\[
-\big(\mathcal{L}^\pi\varphi-\rho_LV-\Lambda\big)(x_L)
= -\tfrac12\sigma^2(x_L,\pi)\,q - \mu(x_L,\pi)p + \rho_LV(x_L)+\Lambda(x_L),
\]
which tends to $-\infty$ as $q\to\infty$ because $\sigma^2(x_L,\pi)>0$.
The supersolution inequality fails, so $J>0$ is impossible.

With Step 1 this forces $J=0$.
\end{proof}

\begin{remark}[SMFN: Why the argument does not prove too much]
\label{rem:smfn}
The same computation run at $J=0$ is consistent, which is the check that
the argument discriminates rather than merely destroys. If $V$ is
differentiable at $x_L$ then any $\varphi$ touching from below must satisfy
$\varphi'(x_L)=1+\kappa$, and since $V$ is affine on the intervention side,
$\varphi\le V$ with matching value and slope forces $\varphi''(x_L)\le0$.
The first entry is then bounded below by
$\rho_LV+\Lambda-(1+\kappa)\mu=\Psi(x_L)$, non-negative by the
intervention-region inequality established in the proof of
Proposition~\ref{prop:soc}(ii) --- a step that itself uses only the
variational inequality, so nothing here is circular. The supersolution property holds. What
breaks the kinked case is not curvature but the freedom a corner creates:
when $V'$ jumps, a test function may match neither one-sided slope, and
that decouples $\varphi''$ from any constraint.

Proposition~\ref{prop:smf} supplies hypothesis (H1) of
Proposition~\ref{prop:ver} at the one point where it was previously
asserted. Interior regularity gives the rest of (H1)
(Lemma~\ref{lem:itr}), and (H2)--(H4) were already established, so
verification is now unconditional.
\end{remark}

\section{Regularity and verification}
\label{sec:verification}

The comparative statics of Section~\ref{sec:compstat-thm} differentiate
$V$, and Lemma~\ref{lem:bdr} assumes joint regularity. This section
records what that regularity amounts to, and what a verification theorem
would and would not require. It closes neither question, but it locates
both more precisely than an assumption does.

\begin{lemma}[ITR: Interior regularity]
\label{lem:itr}
On any compact subset of $(1,\infty)$ the diffusion is uniformly elliptic:
$\sigma^2(x,\pi) \ge (1-|c|)\big[(\pi\sigma x)^2 + \sigma_L^2(1-x)^2\big] > 0$
for $|c|<1$ and $x>1$. On the inaction region the equation is the linear
second-order ODE
$\tfrac12\sigma^2(x,u(x))v'' + \mu(x,u(x))v' - \rho_L v = \Lambda(x)$,
whose coefficients are Lipschitz and are smooth away from the finitely
many points where the cap switches regime or $\Lambda$ kinks. Hence
$v \in C^2$ on the inaction region and $v \in C^3$ away from
$\{R, \bar x\}$, where $\bar x$ is the solvency-liquidity crossing.
\end{lemma}

\begin{proof}
The quadratic form $a^2 + 2cab + b^2$ with $a=\pi\sigma x$ and
$b=\sigma_L(1-x)$ is bounded below by $(1-|c|)(a^2+b^2)$, which is
positive for $x>1$ since then $b \ne 0$. Uniform ellipticity on compacts
plus Lipschitz coefficients gives $v \in C^{2,\alpha}_{\mathrm{loc}}$ by
standard linear theory; where the coefficients are additionally
differentiable, differentiating the equation gives $v'''$ in terms of
lower derivatives, so $v \in C^3$ there. The coefficients fail to be
differentiable only where $u(\cdot)$ switches between the two branches of
the cap, at $\bar x$, and where $\Lambda$ kinks, at $R$.
\end{proof}

\noindent
Two consequences matter for Lemma~\ref{lem:bdr}. The ellipticity
degenerates only at $x=1$, and only because the cap drives $\pi$ to zero
there; the state never reaches that boundary
(Corollary~\ref{cor:boundary}), so the degeneracy is not encountered under
any admissible policy. And the $C^3$ requirement near $y^*$ is met
whenever $y^*$ is separated from $\{R,\bar x\}$, which at the numerical
baseline it is: $R = 1.15$ and $\bar x \approx 1.167$, against
$y^* \approx 1.57$. The $C^3$ requirement is one-sided --- it applies to
the inaction-region solution extended to $y^*$, not across it, since $v$
is linear above $y^*$ and $v'''$ jumps there by construction.

\begin{lemma}[FIN: No chattering]
\label{lem:fin}
For $K>0$, any admissible strategy with finite value makes only finitely
many impulses in expectation, discounted: the discounted count $N$ of
\eqref{eq:Ndef} is finite. Consequently the sum over impulses in the
verification argument converges absolutely.
\end{lemma}

\begin{proof}
Each impulse costs at least $K>0$. A strategy making impulses at times
$\tau_i$ incurs discounted cost at least $K\sum_i e^{-\rho_L\tau_i} = KN$.
The value of any strategy is bounded above by the value of the
unconstrained problem, which is finite under the growth condition below,
so $KN < \infty$ and hence $N<\infty$. The fixed cost, which obstructs
concavity, is what rules out chattering.
\end{proof}

\begin{proposition}[VER: Verification]
\label{prop:ver}
Let $v$ satisfy: (H1) $v \in C^1$ on $(1,\infty)$ and $C^2$ off a finite
set; (H2) the three QVI inequalities of \eqref{eq:qvi} pointwise;
(H3) linear growth, $|v(x)| \le C(1+x)$; and (H4)
$\rho > \sup_x \mu_{\mathrm{opt}}(x)$. Then $v \ge J(\cdot,a)$ for every
admissible $a$, with equality for the barrier-trigger policy; hence
$v = V$ and that policy is optimal.
\end{proposition}

\begin{proof}[Proof sketch]
Apply It\^o to $e^{-\rho_L t}v(X_t)$ under an arbitrary admissible $a$,
splitting the drift, the singular and the impulse parts. The drift term is
non-positive by the first inequality in (H2), since it holds for every
$\pi \le \bar\pi(x)$ and in particular for the one used. The singular term
contributes $\int e^{-\rho_L t} v'(X_t)\,dD_t \ge \int e^{-\rho_L t}dD_t$
by $v' \ge 1$. Each impulse contributes
$v(x+\xi) - v(x) \le K + (1+\kappa)\xi$ by $v \ge \mathcal{M}v$, and the
sum over impulses converges by Lemma~\ref{lem:fin}. Taking expectations
and letting $T\to\infty$, the boundary term
$\mathbb{E}[e^{-\rho_L T}v(X_T)]$ vanishes by (H3) and (H4): the drift is
bounded above by $x(\mu_{\mathrm{opt}}-\mu_L)+\gamma$, so $X$ grows at
most at rate $\sup\mu_{\mathrm{opt}}-\mu_L$, and (H4) is exactly
$\rho_L > \sup\mu_{\mathrm{opt}}-\mu_L$. Reversing the inequalities for
the candidate policy, which attains equality in each, gives the result.
\end{proof}

\begin{remark}[VERN: What verification needs]
\label{rem:vern}
Concavity of $V$ is not among the hypotheses, and its absence
(Remark~\ref{rem:nonconcave}) is therefore not an obstruction to
verification. Concavity is normally used to \emph{derive} the
threshold form of the optimal policy a priori; here that form is imposed
by construction and validated ex post, so the role concavity would play is
already discharged. What is needed in its place is local: $V''(y_{\mathrm{post}})<0$,
so the maximiser in $\mathcal{M}v$ is unique --- a strictly weaker
requirement than global concavity, and the same hypothesis
Proposition~\ref{prop:kcs} already carries. The fixed cost, meanwhile,
cuts both ways: it destroys concavity but supplies
Lemma~\ref{lem:fin}, without which the impulse sum need not converge.

Hypothesis (H4) is the same condition met twice already --- it signs
$dy^*/d\lambda_S$ in Proposition~\ref{prop:tcs}, supplies the
non-degeneracy for Lemma~\ref{lem:bdr}(i), and here delivers
transversality. It is enforced on every solve by the parameter-validity
check of Remark~\ref{rem:solver}.

Hypothesis (H1) is now established rather than assumed. Interior
regularity (Lemma~\ref{lem:itr}) gives $C^2$ and locally $C^3$ away from
the free boundaries; the $C^2$ fit at $y^*$ is the smooth-fit condition
itself; value matching at $x_L$ gives $C^0$; and the $C^1$ fit at $x_L$,
the one point previously asserted, is Proposition~\ref{prop:smf}.
Proposition~\ref{prop:ver} is therefore unconditional.

It is worth recording why the numerics could not settle this and the
argument could. The discrete scheme enforces $v=\max(v,\mathcal{M}v)$
pointwise, so the computed solution carries a corner at whichever node the
boundary falls on, at every mesh; the observed one-sided slope gap is
consistent with a genuine kink and with a steep continuous rise alike, and
separating them requires resolving a layer that the affordable meshes do
not resolve. The supersolution argument avoids the question entirely: it
does not measure the corner, it shows a corner cannot exist.
\end{remark}

\section{The numerical solve}
\label{sec:numerical}

This section records what a numerical solve of \eqref{eq:qvi}
establishes. Its contents are findings at stated parameters, not
theorems. All figures use $r=0.025$, $\mu_s=0.04$, $\mu_L=0.03$,
$\gamma=0.02$, $\rho=0.12$ (so $\rho_L=0.09$), $\sigma=0.08$,
$\sigma_L=0.03$, $c=0.20$,
$a_1=0.045$, $a_2=0.05$, $a_3=0.30$, $K=0.02$, $\kappa=0.01$, $R=1.15$,
on a uniform grid with $y_{\max}=2.5$ and $N=301$ unless stated otherwise.

\begin{remark}[SLV: The solver, and how it was validated]
\label{rem:solver}
The vehicle is the benchmark's own published implementation. Three
validations were run before any conclusion was drawn from it. First,
\emph{reproduction}: at the benchmark's own baseline it returns
$y^*=1.673866$, $y^*_{\mathrm{post}}=1.227682$ and a value uplift of
$0.522875$, against published $1.674$, $1.228$ and $0.5229$ --- four
significant figures on all three. Second, \emph{nesting as a software
test}: $\Lambda$ enters the equation linearly, so it was added as a source
term on the implicit step, and Proposition~\ref{prop:nesting} then
supplies a falsifiable check --- at $\lambda_S=1$ the modified solver must
reproduce the unmodified one \emph{exactly}. It does, to a difference of
zero, and at $\lambda_S=1.5$ the value function is strictly lower below
$R$, confirming the term is active rather than silently dropped.
Third, \emph{operator verification}: the impulse operator $\mathcal{M}$ is
implemented as an $O(N)$ suffix maximum, and is checked against a direct
$O(N^2)$ supremum over the same grid on six profiles --- concave, affine,
non-concave in two distinct shapes, piecewise-flat, and monotone
decreasing --- agreeing to $5.6\times10^{-16}$. The non-concave profiles
are the point of the check: a suffix maximum is most likely to misbehave
exactly where $V$ is not concave, which is the regime
Remark~\ref{rem:nonconcave} reports. Small grids, degenerate cost
parameters, the strictness of $\xi>0$ and the sentinel at $y_{\max}$ are
covered by separate tags.
Convergence was certified by re-solving every $\lambda_S$ at a doubled
iteration cap. All six then converge, at $1645$--$1669$ iterations, and
the induced shifts in $y^*$ fall monotonically from $5.0\times10^{-4}$ to
$4.1\times10^{-4}$ across the range --- $4.9\%$ of the smallest inter-row
gap, so the orderings below survive with about a twentyfold margin. The
recapitalisation boundary is unmoved to machine precision, being pinned to
the mesh. All figures quoted here are from the converged solves.
\end{remark}

\begin{remark}[CBI: The cap binds throughout the inaction region]
\label{rem:capbinds}
Under the solved policy the regulatory constraint is active across the
whole inaction region $[x_L, y^*]$: comparing the solver's chosen
$\pi^\star$ against $\bar\pi$ cell by cell, the two agree at every one of
the $108$--$111$ cells at $\lambda_S \in
\{1.00,\,1.10,\,1.25,\,1.50,\,1.75\}$. At $\lambda_S=2.00$ three cells
show a slack cap, and they straddle the regime switch
$\bar x = 1.167647$, so they are more plausibly an artefact of the switch
than an interior slack region. The comparison is a genuine test rather
than a tautology: the solver searches $\pi$ over a grid spanning
$[0,\bar\pi(x)]$ and retains the maximiser, so a slack optimum would be
found if one existed.

This is a finding at these parameters and not a theorem, and it should be
read at that strength. What it establishes is that the capped
coefficients \eqref{eq:capcoef} describe the whole of the inaction region
here, and that what closes the problem is the impulse control rather than
any interior relaxation of the cap. The formulation \eqref{eq:qvi}
accommodates either outcome without modification.
\end{remark}

\begin{remark}[CSL: Comparative statics in $\lambda_S$]
\label{rem:compstat}
The recapitalisation trigger is monotone in $\lambda_S$, and so is the
dividend barrier. Under the benchmark's proportional-cost operator
$\mathcal{H}$, as $\lambda_S$ runs over
$\{1,\,1.1,\,1.25,\,1.5,\,1.75,\,2\}$: $y^*$ rises
$1.3725 \to 1.4371$, the recapitalisation boundary rises
$1.145 \to 1.200$, and $V(1)$ falls $0.5037 \to 0.4702$. Under the
additive-cost operator $\mathcal{M}$ the same orderings hold and are
larger: $y^*$ rises $1.5715 \to 1.6596$ (a move of $0.088$ against
$0.065$), the boundary rises $1.030 \to 1.105$, and $V(1)$ falls
$0.3843 \to 0.3395$. Every sequence is strictly monotone with no
reversal. The economic reading is that a more asymmetric institution
retains a larger buffer before paying out \emph{and} recapitalises from a
higher level of capital, i.e.\ intervenes earlier; and because the
orderings are common to both cost conventions, they are not an artefact
of either.
\end{remark}

\begin{remark}[LNC: Local non-concavity under a fixed issuance cost]
\label{rem:nonconcave}
Under $\mathcal{M}$ the solve exhibits $V''>0$ on an interval immediately
above the recapitalisation boundary. Three features of it are established
numerically.

\emph{It is not caused by $\Lambda$.} It is present at $\lambda_S=1$,
where $\Lambda\equiv0$ identically and the problem is the risk-neutral
benchmark carrying the additive-cost operator
(Proposition~\ref{prop:nesting}).

\emph{It is not a discretisation artefact.} A kink artefact is pinned to a
fixed number of grid points, so its width in $x$ shrinks with the mesh; a
genuine region has a fixed width, so its point count grows. Over three
meshes at $\lambda_S=1$ the point count went $6 \to 12 \to 24$ --- exactly
the mesh ratio at both refinements --- while the span converged upward
($0.0250$, $0.0275$, $0.0288$). The product of count and spacing is
$0.0300$ at all three meshes, and Richardson extrapolation on the span
gives the same limit, so the region has a converged width of $0.0300$ and
the $N=301$ figure understates it by about a sixth.

\emph{Its mechanism is the fixed cost.} The trigger-to-target distance is
near-constant at $0.365$--$0.370$ across the sweep, and $V'$ must return
to $1+\kappa$ at the target; a fixed cost separates trigger from target
and thereby forces an excursion of $V'$ between them. Local
non-concavity from fixed transaction costs is a familiar feature of
impulse-control value functions, not a pathology peculiar to this model.

The region is present at every $\lambda_S$ tested, occupying $5$ to $8$
grid points at $N=301$ --- a width of $0.025$ to $0.040$, with $V'$
peaking there at $0.121$ to $0.133$ above $1+\kappa$. So $V$ is not
concave for $K>0$ at these parameters. Mesh refinement across $\lambda_S$
has been carried out at $\lambda_S=1$ only.
\end{remark}

\begin{remark}[KSW: The fixed cost is the cause, a $K$-sweep]
\label{rem:Ksens}
Remark~\ref{rem:nonconcave} attributes the convex region to the fixed
issuance cost $K$ rather than to $\Lambda$. That attribution makes a sharp
prediction --- remove $K$ and the region must disappear --- and the
prediction is testable directly. Fixing $\lambda_S=1$, so that
$\Lambda\equiv0$ and $K$ is the only quantity varying, and solving for
$K \in \{0,\,0.001,\,0.005,\,0.01,\,0.02,\,0.05\}$ at $N=301$, every run
converging:

\emph{At $K=0$ everything collapses, exactly.} The trigger and the target
coincide, $V'$ never exceeds $1+\kappa$ anywhere (the largest excess is
$2\times10^{-14}$, machine zero), the convex region has zero width, and
the one-sided derivatives on both sides of the trigger read
$1+\kappa$ to eight decimals. Every feature attributed to the fixed cost
vanishes precisely when the fixed cost vanishes.

\emph{The classical scaling appears unbidden.} The continuation band
$y^*-x_L$ widens monotonically from $0.159$ to $0.685$ across the range,
and the trigger-to-target distance regressed on $K$ in logarithms has
slope $0.3845$, against the $1/3$ that impulse control with a fixed cost
predicts; the ratio of that distance to $K^{1/3}$ sits between $1.26$ and
$1.35$ over the upper four values. Nothing in the construction was fitted
to produce this, so it is an independent check that the solver is doing
impulse control rather than something that merely resembles it.

\emph{The width of the convex region grows with $K$ and then flattens.}
Refining to $N=601$ and extrapolating gives $0.025$ at $K=0.005$ and
$0.030$ at both $K=0.01$ and $K=0.02$, the last agreeing exactly with the
three-mesh value of Remark~\ref{rem:nonconcave}. Widths read off a
single $N=301$ mesh are not converged --- they err in both directions, by
about a sixth at $K=0.02$ and about two fifths at $K=0.005$ --- and an
apparent non-monotonicity in them does not survive refinement.
Single-mesh widths should not be quoted.

At $K=0.05$ the recapitalisation region is one cell wide at $N=301$: the
trigger has collapsed onto the lower boundary, and the diagnostics there
are degenerate rather than informative. That is economically unsurprising
--- a large enough fixed cost makes recapitalisation not worth doing until
ruin is imminent --- but it means the largest $K$ carries no weight in any
of the readings above.
\end{remark}

\section{What is not established}
\label{sec:open}

One thing genuinely open, and one place where the numerics mislead about
something the theory has since settled.

\emph{The relation between $\lambda_V$ and $\lambda_S$.}
Proposition~\ref{prop:statemap} makes this a well-posed question by
putting both on the same state measured from the same origin, and
Proposition~\ref{prop:satlimits} settles that the \emph{definitional}
evaluation of $\lambda_V$ cannot answer it, being independent of
$\lambda_S$. An \emph{estimated} $\lambda_V$ does vary with $\lambda_S$: applying the
panel estimator to simulated model paths (quarterly sampling, sixty
quarters, four hundred paths per parameter value) gives a median that
falls monotonically from about $1.33$--$1.39$ at $\lambda_S=1$ to about
$1.15$--$1.17$ at $\lambda_S=2$, depending on the issuance-cost
convention. The direction is known, but the movement is small relative to
the noise it is measured against: roughly $0.18$--$0.22$ of total change
across the full admissible range, against a per-path standard deviation
of $0.16$--$0.20$. The estimated ratio is therefore not uninformative but
weakly identified --- distinguishable from noise only when aggregated
across many institutions, not from any single one. Any empirical reading
that compares an estimate of $\lambda_V$ against a benchmark value should
report both the direction above and this signal-to-noise magnitude,
rather than treat a single point estimate as diagnostic on its own.

\emph{Regularity at the recapitalisation trigger, and why the numerics
could not have shown it.} That $V$ is $C^1$ at the trigger is settled by
Proposition~\ref{prop:smf}, but the route is worth recording because the
solve is misleading on the point. For $K>0$ the one-sided derivatives at
the detected trigger differ in every solve, which would appear to say $V$
is $C^0$ but not $C^1$ there. That reading is an artefact of the scheme:
the detected trigger is the last mesh point at which the impulse obstacle
is active, so the true boundary lies inside the following cell, and a
one-sided difference across that cell mixes the two branches in proportion
to where the boundary falls; that proportion moves essentially at random
as parameters shift the boundary within a cell, and at fixed mesh and
fixed $K$ the mixed quantity varies by a factor of sixteen across
$\lambda_S$, non-monotonically, which no genuine property of $V$ would do.
Mesh refinement does not escape this: the scheme enforces
$v=\max(v,\mathcal{M}v)$ pointwise, so a corner sits at whichever node the
boundary occupies at every resolution. The question is therefore not one
the mesh can answer, and Proposition~\ref{prop:smf} answers it instead ---
not by measuring the corner but by showing, from the variational
inequality and the supersolution property, that a corner cannot exist.

\appendix
\section{Verification status}
\label{app:status}

Each result has one canonical verifier: a Lean proof where one exists,
an exact symbolic check otherwise. Every Lean declaration cited below is
compiler-verified and depends on no \texttt{sorry} and no axiom beyond
Lean's classical defaults. Four results are marked \emph{Lean (partial)}:
the algebraic core of the argument is compiler-verified with the analytic
input (an envelope identity, a cost bound) stated as an explicit
hypothesis, and the row records which step remains proved only in the
text. Symbolic checks terminate in
\texttt{simplify(\dots)==0} or a sign condition on symbols, with no
floating point and no sampling. Numerical findings are labelled as such;
a numeric check can only fail to find a counterexample and is evidence,
not proof.

\begin{longtable}{p{0.34\textwidth} p{0.16\textwidth} p{0.42\textwidth}}
\toprule
\textbf{Result} & \textbf{Verifier} & \textbf{Location} \\
\midrule
\endhead
Prop.~\ref{prop:identity}--\ref{prop:persistence} & Lean &
\texttt{RateBased.lean} \\
Thm.~\ref{thm:lambda4} & symbolic &
\texttt{verify\_rate\_based.py} (S1a--S1b) \\
Prop.~\ref{prop:robust}, \ref{prop:egarch}, \ref{prop:sojourn} &
symbolic & \texttt{verify\_rate\_based.py} \\
Prop.~\ref{prop:nesting} & Lean &
\texttt{QVI\_Part1.lean}, \texttt{lambda\_asymmetry\_vanishes\_at\_one} \\
Prop.~\ref{prop:concave-payoff} & Lean &
\texttt{QVI\_Part1.lean}, \texttt{neg\_Lambda\_concave} \\
Prop.~\ref{prop:rrL} & symbolic & \texttt{verify\_state\_map.py} (M10) \\
Prop.~\ref{prop:lowerbound} & proof in text &
admissibility construction \\
Prop.~\ref{prop:statemap}, Cor.~\ref{cor:boundary} & symbolic &
\texttt{verify\_state\_map.py} (M1--M9) \\
Prop.~\ref{prop:satlimits} & symbolic &
\texttt{verify\_saturation\_limits.py} (L1--L3; L5* expected-fail) \\
Lem.~\ref{lem:reg} & proof in text &
convexity of the value in each parameter; differentiable off a countable
set \\
Lem.~\ref{lem:envelope} & Lean (partial) &
\texttt{Envelope.lean}, \texttt{wode\_of\_envelope}: the ODE
\eqref{eq:Wode} from the envelope identity \eqref{eq:envelope} as
hypothesis; the Danskin step proving \eqref{eq:envelope} is in the text \\
Lem.~\ref{lem:bdr} & proof in text &
implicit function theorem at each boundary \\
Prop.~\ref{prop:tcs} & symbolic &
\texttt{verify\_comparative\_statics.py} (H4a, H4b, H4d) \\
Rem.~\ref{rem:tcsn} & numerical &
\texttt{verify\_document\_figures.py} (F5) \\
Lem.~\ref{lem:envelopeK} & Lean (partial) &
\texttt{Envelope.lean}, \texttt{node\_of\_envelope},
\texttt{neumann\_at\_ystar}, \texttt{trigger\_jump}: \eqref{eq:Node}
and both boundary data from \eqref{eq:envelopeK}, the first-order
condition, smooth fit and value matching as hypotheses; the Danskin step
proving \eqref{eq:envelopeK} is in the text \\
Prop.~\ref{prop:kcs} & symbolic &
\texttt{verify\_comparative\_statics.py} (H4c, H4e) \\
Cor.~\ref{cor:idn} & symbolic &
\texttt{verify\_comparative\_statics.py} (H5a--H5d) \\
Rem.~\ref{rem:kcsn} & numerical &
six $K$ values; directions at stated parameters \\
Prop.~\ref{prop:soc} & symbolic &
\texttt{verify\_comparative\_statics.py} (H3a) \\
Rem.~\ref{rem:socn} & numerical &
\texttt{verify\_document\_figures.py} (F6--F8) \\
Prop.~\ref{prop:smf} & Lean &
\texttt{SmoothFit.lean}, \texttt{smooth\_fit} \\
Rem.~\ref{rem:smfn} & Lean (partial) &
\texttt{Envelope.lean}, \texttt{smf\_permits}: every curvature
$q \le 0$ satisfies the operator inequality given
$\Psi(x_L) \ge 0$; that these are exactly the admissible curvatures
under smooth fit is argued in the text \\
Lem.~\ref{lem:itr} & proof in text &
interior trigger, $x_L > 1$ \\
Lem.~\ref{lem:fin} & Lean (partial) &
\texttt{ImpulseCount.lean}, \texttt{count\_summable},
\texttt{count\_le}, \texttt{impulse\_sum\_summable}: summability of the
discounted count with $N \le B/K$, and absolute convergence of the
impulse sum, from the cost bound as hypothesis; that each impulse costs
at least $K$ under finite value is in the text \\
Prop.~\ref{prop:ver} & proof in text &
It\^o argument on $e^{-\rho_L t}v(X_t)$; hypotheses (H1)--(H4) \\
Rem.~\ref{rem:vern} & numerical &
\texttt{verify\_smooth\_fit\_shooting.py} (S2--S3) \\
Rem.~\ref{rem:solver} & numerical &
benchmark reproduced to 4 s.f.; $\lambda_S=1$ nesting exact;
$\mathcal{M}$ vs direct supremum, \texttt{verify\_M\_operator.m}
(V1--V10); the $\mathcal{M}$ sweep convergence-certified \\
Rem.~\ref{rem:capbinds} & numerical &
$\pi^\star$ vs $\bar\pi$ across the inaction region, six $\lambda_S$ \\
Rem.~\ref{rem:compstat} & numerical &
six $\lambda_S$, both operators, strict monotonicity, no reversal; the
$\mathcal{H}$ figures are from a sweep not convergence-certified \\
Rem.~\ref{rem:nonconcave} & numerical &
mesh-refinement discriminant over three meshes, width converged to
$0.0300$; present at all six $\lambda_S$ \\
Rem.~\ref{rem:Ksens} & numerical &
six $K$ values, all converged; $K=0$ control exact to machine precision;
$K^{1/3}$ scaling recovered independently \\
\bottomrule
\end{longtable}

\end{document}
